\documentclass[11pt, a4paper]{article}

% --- Packages ---
\usepackage[utf8]{inputenc}
\usepackage{graphicx}
\usepackage[margin=0.75in]{geometry}
\usepackage{times}
\usepackage{xcolor}
\usepackage{titlesec}
\usepackage{float}
\usepackage{booktabs} 
\usepackage{multirow}
\usepackage{tabularx}
\usepackage{array} 
\usepackage{enumitem}
\usepackage{pifont} 
\newcommand{\tick}{\ding{52}} 
\newcommand{\cross}{\ding{56}} 
\usepackage{caption}
\usepackage{amsmath}
\usepackage{amssymb}

% --- Amrita Colors ---
\definecolor{amritaMaroon}{HTML}{A4123F}

% --- Title & Section Formatting ---
\titlespacing*{\section}{0pt}{12pt}{4pt}
\titleformat{\section}{\large\bfseries\color{amritaMaroon}}{\thesection}{1em}{}
\titlespacing*{\subsection}{0pt}{9pt}{3pt}
\titleformat{\subsection}{\normalsize\bfseries}{\thesubsection}{1em}{}
\titlespacing*{\subsubsection}{0pt}{6pt}{2pt}
\titleformat{\subsubsection}{\small\bfseries}{\thesubsubsection}{1em}{}

\pagestyle{empty}

\begin{document}

% ==========================================
% HEADER
% ==========================================
\noindent
\begin{minipage}{\textwidth}
    \centering
    {\Large \textbf{\textcolor{amritaMaroon}{
        AMRITA VISHWA VIDYAPEETHAM \\
        SCHOOL OF COMPUTING \\
        DEPARTMENT OF COMPUTER SCIENCE AND ENGINEERING
    }}}
    \vspace{0.2cm} \\
    {\large \textbf{AI-Driven Forensic Indexing: A Behavioral Model for Chain Snatching Identification}} \\
    \vspace{0.1cm}
    23CSE399 -- Project Phase 2 (Panel Review 1) \\ \today
\end{minipage}

\vspace{0.3cm}
\hrule
\vspace{0.3cm}

% --- Team & Guide Details ---
\noindent
\begin{center}
\renewcommand{\arraystretch}{1.4}
\setlength{\tabcolsep}{5pt}
\begin{tabular}{|l|p{4.2cm}|p{4cm}|p{2.5cm}|c|}
\hline
\multicolumn{2}{|c|}{\textbf{Team Members (Team 24)}} & \textbf{Guide} & \textbf{Project Type} & \textbf{Marks} \\
\hline
\small CB.SC.U4CSE23151 & \small Vasudev Kishor &
    \multirow{4}{=}{\raggedright \small \textbf{Ms.~Sathiya R R} \newline \textit{Assistant Professor}} &
    \multirow{4}{=}{\raggedright \small \textbf{Type 1} \newline Research} &
    \multirow{4}{*}{\textbf{\huge \dots\ / 20}} \\
\cline{1-2}
\small CB.SC.U4CSE23342 & \small P Dharshni & & & \\
\cline{1-2}
\small CB.SC.U4CSE23710 & \small Rathish B & & & \\
\cline{1-2}
\small CB.SC.U4CSE23750 & \small Thilagan Iniyavan & & & \\
\hline
\end{tabular}
\end{center}

\section{Problem Statement \& Research Context}

\subsection{Problem Statement}
Municipal CCTV surveillance networks generate vast streams of high-resolution video continuous footage. In criminal forensics—specifically for micro-duration events like chain snatching—investigators face severe processing bottlenecks. Chain-snatching events usually last only a few seconds within hours of uninformative footage, rendering manual inspection unsustainable.

Standard computer vision approaches face significant operational limitations:
\begin{itemize}[leftmargin=*, noitemsep]
    \item \textbf{Manual Sifting Bottleneck:} Law enforcement spends hundreds of hours scanning static or non-event footage.
    \item \textbf{Semantic Intent Deficit:} Bounding-box object detectors (e.g., standard YOLO) locate pedestrians and motorcycles but cannot discern malicious intent or coordinate interaction patterns.
    \item \textbf{Temporal Spread of Evidence:} Criminal signatures are non-instantaneous; malicious intent is distributed chronologically across approach, interaction, snatching, and escape phases.
    \item \textbf{Environmental Degradation:} Surveillance feeds suffer from night-time motion blur, heavy occlusions, 2D projection overlap, and small distant target resolution.
\end{itemize}

\subsection{Research Gap Analysis}
\begin{table}[H]
\centering
\small
\renewcommand{\arraystretch}{1.3}
\begin{tabularx}{\textwidth}{|>{\hsize=0.25\hsize\bfseries}X|>{\hsize=0.35\hsize}X|>{\hsize=0.4\hsize}X|}
\hline
\textbf{Current Method} & \textbf{Capability Provided} & \textbf{Identified Forensic Gap} \\
\hline
Object Detection & Pinpoints spatial location of targets. & Lacks temporal interaction context; high false positive rate for normal riding. \\
\hline
Object Tracking & Maintains target identity across frames. & Fails to recognize coordinate behavioral intent or kinetic theft patterns. \\
\hline
Pose / Action Models & Extracts detailed joint kinematics. & Extremely computationally expensive; unfeasible to run continuously across all video streams. \\
\hline
End-to-End Classifiers & Outputs direct video-level predictions. & Black-box architecture; provides no verifiable chain of evidence or component traceability. \\
\hline
Manual Indexing & Provides human verification. & Unscalable to large-scale municipal CCTV archives. \\
\hline
\end{tabularx}
\caption{Systematic gap analysis of existing forensic video processing approaches.}
\end{table}

\section{Literature Survey \& Model Selection Strategy}

To establish an optimal pipeline architecture, extensive comparative evaluation of state-of-the-art algorithms was conducted across preprocessing, object detection, and action recognition domains.

\subsection{Comparative Model Analysis}

\textbf{A. Low-Light Enhancement Methods}
\begin{table}[H]
\centering
\scriptsize
\renewcommand{\arraystretch}{1.2}
\begin{tabularx}{\textwidth}{|l|X|X|c|}
\hline
\textbf{Model} & \textbf{Core Technique} & \textbf{Primary Limitation} & \textbf{Selection Status} \\
\hline
LIME & Illumination map estimation via structure-aware smoothing & High per-frame latency; over-exposure in dynamic areas & Low Suitability \\
\hline
EnlightenGAN & Unpaired image-to-image translation (GAN) & Risk of artifact generation and feature hallucination & Rejected \\
\hline
Zero-DCE & Non-generative dynamic range curve estimation & Requires scene parameter tuning & \textbf{Selected Engine} \\
\hline
\end{tabularx}
\end{table}

\textbf{B. Object Detection Architectures}
\begin{table}[H]
\centering
\scriptsize
\renewcommand{\arraystretch}{1.2}
\begin{tabularx}{\textwidth}{|l|c|c|c|X|}
\hline
\textbf{Model} & \textbf{mAP@50-95} & \textbf{Latency (ms)} & \textbf{Params} & \textbf{Small Object / Occlusion Handling} \\
\hline
YOLOv8 (Nano) & 37.3\% & 12--15 ms & 3.2M & Moderate baseline accuracy. \\
\hline
YOLOv10 (Nano) & 38.5\% & 10--12 ms & 2.3M & High performance via NMS-free dual assignments. \\
\hline
YOLOv11 (Nano) & 39.5\% & 8--10 ms & 2.6M & Ultra-high accuracy with C2PSA attention modules. \\
\hline
\end{tabularx}
\end{table}

\textbf{C. Action Recognition \& Behavioral Models}
\begin{table}[H]
\centering
\scriptsize
\renewcommand{\arraystretch}{1.2}
\begin{tabularx}{\textwidth}{|X|c|c|X|c|}
\hline
\textbf{Architecture} & \textbf{Accuracy (NTU-120)} & \textbf{GFLOPs} & \textbf{Night \& Occlusion Robustness} & \textbf{Pipeline Selection} \\
\hline
ViViT & 89.1\% & 50.0+ & Low (High pixel dependency) & Rejected (Too heavy) \\
\hline
2s-AGCN & 91.2\% & 18.4 & High (Skeleton-based spatial graph) & Candidate \\
\hline
ST-GCN & 89.9\% & 16.2 & High (Lightweight Spatio-Temporal Graph) & \textbf{Selected Engine} \\
\hline
\end{tabularx}
\caption{Literature benchmarks driving baseline component choices.}
\end{table}

\section{System Architecture \& Mathematical Foundations}

\begin{center}
\IfFileExists{methodology.png}{%
\includegraphics[width=0.92\textwidth,height=0.32\textheight,keepaspectratio]{methodology.png}\\
}{%
\fbox{\parbox[c][3.5cm][0.9\textwidth]{\centering Methodology diagram placeholder: \texttt{methodology.png} not found.}}\\
}
{\scriptsize Figure 1: Proposed Progressive Forensic Indexing Architecture}
\end{center}

\subsection{Architectural Principle \& 17-Stage Pipeline}
The architecture enforces a cascade principle: inexpensive filters process the raw video stream, passing candidate regions to increasingly sophisticated neural networks only when specific criteria are satisfied.

\begin{table}[H]
\centering
\small
\renewcommand{\arraystretch}{1.25}
\begin{tabularx}{\textwidth}{|c|>{\hsize=0.32\hsize}X|>{\hsize=0.68\hsize}X|}
\hline
\textbf{Stage} & \textbf{Main Operational Question} & \textbf{Module Output} \\
\hline
1 & Is there motion? & Motion triage candidate frames. \\
\hline
2--3 & Who is present in the scene? & Object tracks for person and vehicle entities. \\
\hline
4--7 & Are entities interacting spatially? & Spatial relationship vectors and interaction pairs. \\
\hline
8--10 & What behavior trajectory is occurring? & Behavior graphs and localized Region of Interest (ROI) crops. \\
\hline
11--13 & What human action is occurring? & Extracted skeleton keypoints and ST-GCN action labels. \\
\hline
14--15 & Does multi-modal evidence agree? & Fused interaction probability score. \\
\hline
16 & Does it match a snatch signature? & Verified signature score against chronological theft rules. \\
\hline
17 & How is it forensically retrieved? & Indexed event records linked to evidence thumbnails and video clips. \\
\hline
\end{tabularx}
\caption{The 17-stage progressive indexing pipeline operational abstraction.}
\end{table}

\subsection{Mathematical Foundations}

\textbf{1. Zero-DCE Illumination Enhancement Curve:}
To stabilize feature extraction under low-light night-time conditions without adding generative overhead:
\begin{equation}
L_n(x) = L_{n-1}(x) + \alpha_n L_{n-1}(x)(1 - L_{n-1}(x))
\end{equation}
where $L_n(x)$ is the enhanced image channel at iteration $n$, and $\alpha_n \in [-1, 1]$ represents trainable curve parameter maps.

\textbf{2. 3D Spatial Proximity & Perspective Correction:}
To suppress false positive interactions caused by 2D camera projection overlaps, distance $d$ between pedestrian target $s$ and vehicle target $v$ is calculated with depth-scaling factors:
\begin{equation}
d = \sqrt{w_x(x_s - x_v)^2 + w_y(y_s - y_v)^2 + w_z(z_s - z_v)^2}
\end{equation}

\textbf{3. Spatio-Temporal Graph Convolution (ST-GCN):}
Given a joint skeleton graph $G = (V, E)$ with joint nodes $V$ and bone/temporal edges $E$, layer-wise spatial graph convolutions are computed as:
\begin{equation}
f_{out} = \sum_{k}^{K_v} \Lambda_k^{-\frac{1}{2}} A_k \Lambda_k^{-\frac{1}{2}} f_{in} W_k
\end{equation}
where $A_k$ represents the adjacency matrix for spatial partition subset $k$, $\Lambda_k$ is the normalized diagonal degree matrix, and $W_k$ represents the trainable weight vector.

\section{Core Intelligent Modules}

The intelligent processing capabilities of the platform are organized into four primary modules.

\subsection{Behavior Intelligence \& Relationship Engine}
\begin{table}[H]
\centering
\small
\renewcommand{\arraystretch}{1.3}
\begin{tabularx}{\textwidth}{|>{\hsize=0.3\hsize\bfseries}X|>{\hsize=0.7\hsize}X|}
\hline
\textbf{Sub-Module} & \textbf{Functional Operation \& Engineering Design} \\
\hline
Motion Feature Engine & Derives kinematics including instant velocity ($\vec{v}$), acceleration ($\vec{a}$), directional vectors, and trajectory history. \\
\hline
Relationship Engine & Evaluates dynamic bounding box spatial overlaps, relative approach rates, and bounding box scale changes between entities. \\
\hline
Interaction Manager & Tracks entity-pair interactions over time using sliding window buffers, identifying sustained convergence. \\
\hline
Behavior Graph Builder & Encapsulates interaction states into Directed Acyclic Graphs (DAGs) representing temporal transitions (Approach $\rightarrow$ Proximity $\rightarrow$ Separation). \\
\hline
\end{tabularx}
\caption{Behavior Intelligence Layer design specifications.}
\end{table}

\subsection{Kinematic Pose \& ST-GCN Action Recognition Engine}
\begin{table}[H]
\centering
\small
\renewcommand{\arraystretch}{1.3}
\begin{tabularx}{\textwidth}{|>{\hsize=0.3\hsize\bfseries}X|>{\hsize=0.7\hsize}X|}
\hline
\textbf{Feature Dimension} & \textbf{Engine Architecture \& Specification} \\
\hline
Input ROI Selection & Crop regions strictly isolated to interacting person--vehicle bounding boxes. \\
\hline
Pose Extraction & Baseline MediaPipe engine generating 2D joint coordinate matrices ($17 \text{ keypoints} \times \text{confidence}$). \\
\hline
Spatio-Temporal Graph & Maps joints as graph vertices ($V$) across frame temporal sequences ($T$). \\
\hline
Action Output & Predicts action probabilities for critical snatch primitives: \textit{Reaching}, \textit{Grabbing}, \textit{Pulling}, and \textit{Normal Riding}. \\
\hline
\end{tabularx}
\caption{Pose Estimation and ST-GCN Module specifications.}
\end{table}

\subsection{Snatch Signature Verification Layer}
\begin{table}[H]
\centering
\small
\renewcommand{\arraystretch}{1.3}
\begin{tabularx}{\textwidth}{|c|>{\hsize=0.35\hsize\bfseries}X|>{\hsize=0.55\hsize}X|}
\hline
\textbf{Sequence Step} & \textbf{Behavioral Primitive} & \textbf{Required Evidence Metric} \\
\hline
Step 1 & Approach / Convergence & Decreasing spatial distance trajectory ($d/dt < 0$). \\
\hline
Step 2 & Proximity Interaction & Intersection over Union (IoU) / Proximity threshold hit. \\
\hline
Step 3 & Upper-Limb Theft Action & High-confidence ST-GCN detection (\textit{Reach/Grab/Pull}). \\
\hline
Step 4 & Rapid Acceleration & High positive velocity gradient ($d\vec{v}/dt \gg 0$). \\
\hline
Step 5 & Escape Divergence & Rapid spatial divergence ($d/dt \gg 0$). \\
\hline
\end{tabularx}
\caption{Chronological theft signature matching verification steps.}
\end{table}

\subsection{Forensic Metadata Indexing Module}
\begin{table}[H]
\centering
\small
\renewcommand{\arraystretch}{1.3}
\begin{tabularx}{\textwidth}{|>{\hsize=0.3\hsize\bfseries}X|>{\hsize=0.7\hsize}X|}
\hline
\textbf{Index Field} & \textbf{Description and Query Utility} \\
\hline
Temporal Metadata & Precise start/end frame indices, system timestamps, and video ID. \\
\hline
Entity Association & Track ID links corresponding to involved pedestrian and vehicle instances. \\
\hline
Behavioral Labels & Graph transition logs, action predictions, and signature confidence scores. \\
\hline
Evidence References & Direct memory offsets to saved keyframe thumbnails and short evidence video clips. \\
\hline
Search Capabilities & Supports instant SQL-like query execution (e.g., \texttt{WHERE Action='Grabbing' AND Confidence > 0.85}). \\
\hline
\end{tabularx}
\caption{Data schema for the forensic event database.}
\end{table}

\section{Research Strategy \& Experimental Matrix}

\subsection{Research Strategy \& Integrity Controls}
To maintain rigorous scientific standards as a Type 1 Research Project, experiments enforce strict protocol controls:

\begin{table}[H]
\centering
\small
\renewcommand{\arraystretch}{1.3}
\begin{tabularx}{\textwidth}{|>{\hsize=0.35\hsize\bfseries}X|>{\hsize=0.65\hsize}X|}
\hline
\textbf{Research Challenge} & \textbf{Integrity Control Protocol} \\
\hline
Benchmarking Bias & All model comparisons are evaluated on identical, standardized CCTV video benchmarks. \\
\hline
Configuration Drift & Hyperparameters, hardware constraints, and seeds are auto-logged per test run. \\
\hline
Missing Implementations & Unavailable external models are recorded strictly as \textit{not evaluated} (zero artificial approximation). \\
\hline
System Overhead Analysis & Automated tracking of CPU memory, GPU VRAM allocations, and frame-rate throughput. \\
\hline
\end{tabularx}
\caption{Protocol controls for experimental validity.}
\end{table}

\subsection{Comprehensive Experimental Matrix}
\begin{table}[H]
\centering
\scriptsize
\renewcommand{\arraystretch}{1.25}
\begin{tabularx}{\textwidth}{|>{\hsize=0.18\hsize}X|>{\hsize=0.32\hsize}X|>{\hsize=0.32\hsize}X|>{\hsize=0.18\hsize}X|}
\hline
\textbf{Pipeline Stage} & \textbf{Candidate Methods} & \textbf{Evaluation Metrics} & \textbf{Status} \\
\hline
Motion Triage & Frame Difference, MOG2, KNN & Reduction Rate, Recall, FPS & \textbf{Completed} \\
\hline
Object Detection & YOLOv8, YOLOv10, YOLOv11 & mAP@50, Precision, Latency & Planned \\
\hline
Tracking & ByteTrack, BoT-SORT & ID Switches, MOTA, Continuity & Planned \\
\hline
Relationship & Proximity, Overlap, Scale Change & Precision, Recall & Planned \\
\hline
Pose Extraction & MediaPipe, AlphaPose & Keypoint PCK, Latency & \textbf{MediaPipe Active} \\
\hline
Action Model & ST-GCN vs Executable Alts & F1-Score, Processing Latency & \textbf{ST-GCN Ready} \\
\hline
Evidence Fusion & Weighted, Rule, Bayesian, Voting & Agreement Rate, F1-Score & \textbf{Implemented} \\
\hline
Architecture & Baseline vs Progressive Cascade & Throughput (FPS), End-to-End F1 & \textbf{Implemented} \\
\hline
Forensic Indexing & Inverted Index vs Linear Scan & Query Latency vs Record Volume & Synthetic Active \\
\hline
\end{tabularx}
\caption{Master experimental matrix for model evaluation and component ablation.}
\end{table}

\section{Summary of Completed Work ($>$60\% Milestone)}

The fundamental system framework is fully operational, boasting an active verification suite with \textbf{164 out of 164 unit tests passing}.

\subsection{Completed Study: Motion Triage Performance}
An empirical study was conducted to establish the optimal motion filtering setup to maximize frame reduction while retaining subtle event triggers.

\begin{table}[H]
\centering
\small
\renewcommand{\arraystretch}{1.3}
\begin{tabularx}{\textwidth}{|>{\hsize=0.3\hsize\bfseries}X|>{\hsize=0.7\hsize}X|}
\hline
\textbf{Parameter} & \textbf{Empirical Result / Observation} \\
\hline
Evaluated Methods & Frame Differencing, Mixture of Gaussians (MOG2), K-Nearest Neighbors (KNN). \\
\hline
Selected Algorithm & \textbf{Frame Differencing} (Optimal balance between compute overhead and recall). \\
\hline
Selected Threshold & \textbf{2000 pixels (Setting 20)}. \\
\hline
Processing Speed & \textbf{$\approx 335$ FPS} on standard benchmark footage. \\
\hline
Sensitivity Trade-off & Lower thresholds ($<1000$ px) retain noise; higher thresholds ($>3000$ px) risk missing subtle motion buildups. \\
\hline
\end{tabularx}
\caption{Empirical results for Stage 1 Motion Triage.}
\end{table}


\subsection{Current Implementation Progress}

\begin{table}[H]
\centering
\small
\renewcommand{\arraystretch}{1.3}
\begin{tabularx}{\textwidth}{|>{\hsize=0.4\hsize}X|c|>{\hsize=0.6\hsize}X|}
\hline
\textbf{Module} & \textbf{Status} & \textbf{Current evidence} \\
\hline
Motion triage & Completed & Comparison study \\
\hline
Detection and tracking & Completed & Integrated pipeline \\
\hline
Interaction and behaviour & Completed & Integration tests \\
\hline
Behaviour graph & Completed & Integration tests \\
\hline
ROI and skeleton preparation & Completed & Pipeline tests \\
\hline
Pose estimation & Implemented & MediaPipe baseline \\
\hline
ST-GCN action recognition & Implemented & Integrated stage \\
\hline
Behaviour fusion & Completed & Multiple strategies \\
\hline
Snatch signature engine & Completed & Signature evaluation \\
\hline
Forensic indexing & Completed & Query engine \\
\hline
Research comparison framework & Implemented & Comparison / ablation runners \\
\hline
\end{tabularx}
\end{table}

\begin{center}
\textbf{164/164 unit tests passing in the current implementation.}
\end{center}


\subsection{Current Research Findings}
\begin{table}[H]
\centering
\small
\renewcommand{\arraystretch}{1.3}
\begin{tabularx}{\textwidth}{|>{\hsize=0.4\hsize}X|>{\hsize=0.6\hsize}X|}
\hline
\textbf{Research Observation} & \textbf{Operational Takeaway} \\
\hline
Frame-level motion triage discards non-events & Bypasses expensive neural network processing on static frames, resolving power bottlenecks. \\
\hline
Spatial overlap triggers false alarms & Spatial proximity alone is insufficient; chronological behavioral context is mandatory. \\
\hline
Behavior Graphs preserve temporal order & Enables representing theft as a dynamic state sequence rather than isolated spatial frame detections. \\
\hline
Skeleton-based action models add robustness & Human joint dynamics effectively complement vehicle trajectory tracking in low-light/occluded scenes. \\
\hline
Forensic metadata query eliminates rescanning & Investigators search timestamped structured records directly, achieving near-instant clip retrieval. \\
\hline
\end{tabularx}
\caption{Key findings derived from initial system testing.}
\end{table}

\section{Plan for Upcoming Semesters}

With core modules active and unit-tested, work in upcoming semesters will execute the comparative research roadmap:

\begin{table}[H]
\centering
\small
\renewcommand{\arraystretch}{1.3}
\begin{tabularx}{\textwidth}{|c|>{\hsize=0.3\hsize\bfseries}X|>{\hsize=0.6\hsize}X|}
\hline
\textbf{Phase} & \textbf{Target Task} & \textbf{Planned Deliverables \& Metrics} \\
\hline
Phase 1 & Detector \& Tracker Benchmarks & Evaluate YOLOv8, v10, v11 and ByteTrack vs. BoT-SORT on standard CCTV datasets for precision and ID continuity. \\
\hline
Phase 2 & Action Model Benchmarks & Benchmark executable pose/action models against ST-GCN baseline for F1-score and frame throughput. \\
Phase 3 & Full Pipeline Ablation & Systematically deactivate individual modules (e.g., ST-GCN, Relationship Engine) to quantify accuracy vs. FPS trade-offs. \\
\hline
Phase 4 & Edge Case & Study failure modes including severe occlusion, low resolution, and distant camera perspectives. \\
\hline
Phase 5 & Final Dissertation & Measure forensic query scaling latency and publish final project research findings. \\
\hline
\end{tabularx}
\caption{Detailed execution plan for Phase 2 implementation.}
\end{table}

\nocite{*}
\bibliographystyle{unsrt}
\bibliography{reff}

\end{document}